\section{A collision-rate obstruction to spatial Laguerre comparisons}
\label{sec:collisionrate}
The heat hierarchy \eqref{eq:backwardlaguerre} suggests a
spatially dependent comparison
\begin{equation}\label{eq:spatialcomparison}
k_n U_{n+1}(s,x)\le c_n(s,x)U_n(s,x),
\qquad k_n=(n+1)(2n+1).
\end{equation}
The bounded-coefficient version was already shown inadequate
for general positive Fourier kernels. Allowing dependence on
$x$ does not automatically cure the problem: a zero collision
forces a nonintegrable singularity in the coefficient.

\begin{theorem}[Universal collision-rate obstruction]
\label{thm:collisionrate}
Suppose a real entire family $H_t$ is smooth in $t$,
satisfies $\partial_tH_t=-\partial_x^2H_t$, and has a
nondegenerate real double zero at $(t_*,x_*)$:
\[
H_{t_*}(x_*)=\partial_xH_{t_*}(x_*)=0,
\qquad a:=\partial_x^2H_{t_*}(x_*)\ne0.
\]
Write $s_*=\tfrac12-t_*$ and
$U_n(s,x)=L_{n,1/2-s}(x)$. Then for $s<s_*$
sufficiently close to $s_*$, $U_1(s,x_*)>0$,
and
\begin{equation}\label{eq:collrate}
\boxed{\lim_{s\uparrow s_*}
(s_*-s)\frac{6U_2(s,x_*)}{U_1(s,x_*)}
=\frac32.}
\end{equation}
Any nonnegative $c_1$ satisfying \eqref{eq:spatialcomparison}
throughout that interval must therefore obey
\[
\liminf_{s\uparrow s_*}(s_*-s)c_1(s,x_*)
\ge\frac32,\qquad
\int_{s_*-\eta}^{s_*}c_1(s,x_*)\,ds=+\infty
\]
for every sufficiently small $\eta>0$.
\end{theorem}
\begin{proof}
Set $\delta=t-t_*>0$. The backward heat equation
and Taylor's theorem at the fixed spatial coordinate give
\[
\begin{aligned}
H_{t_*+\delta}(x_*)&=-a\delta+O(\delta^2),&
\partial_xH_{t_*+\delta}(x_*)&=O(\delta),\\
\partial_x^2H_{t_*+\delta}(x_*)&=a+O(\delta).
\end{aligned}
\]
Consequently
$L_{1,t_*+\delta}(x_*)=a^2\delta+O(\delta^2)$.
The exact coefficient formula
\[
L_{2,t}
=\tfrac1{12}HH^{(4)}
-\tfrac13H'H^{(3)}
+\tfrac14(H'')^2
\]
gives $L_{2,t_*+\delta}(x_*)=a^2/4+O(\delta)$.
Division yields \eqref{eq:collrate}.
Since $c_1\ge6U_2/U_1$, it follows that
$c_1(s,x_*)\ge(3/2+o(1))/(s_*-s)$.
Integrating this necessary lower bound proves divergence.
\end{proof}

\begin{proposition}[Exact two-root gap model]
\label{prop:collisionpair}
The polynomial heat solution
\[
H_t(x)=(x-x_*)^2-2(t-t_*)
\]
has
\[
L_{1,t}(x)=2(x-x_*)^2+4(t-t_*),
\qquad L_{2,t}(x)=1.
\]
For $t>t_*$ the exact adjacent-order quotient is
\begin{equation}\label{eq:pairquotient}
\boxed{\frac{6L_{2,t}(x)}{L_{1,t}(x)}
=\frac{3}{(x-x_*)^2+2(t-t_*)}.}
\end{equation}
The two real zeros at this time have gap
$g=2\sqrt{2(t-t_*)}$, and at their midpoint
every admissible coefficient in \eqref{eq:spatialcomparison}
for $n=1$ must satisfy $c_1\ge12/g^2$.
\end{proposition}
\begin{proof}
All claims follow by direct differentiation.
The two roots are $x_*\pm\sqrt{2(t-t_*)}$,
so $g^2=8(t-t_*)$ and
$3/[2(t-t_*)]=12/g^2$.
\end{proof}

The polynomial model supplies a normal form, not a collision
theorem for the Riemann $\Xi$ function. It clarifies the
proposed logarithmic-square spatial majorant: such a bound
would need control of small zero gaps, as well as collisions,
in the particular heat family being studied. The positive
Gaussian--cosh example of
Proposition~\ref{prop:doublezerofailure} supplies an
actual Fourier-transform heat flow in which a collision
occurs at a finite positive time. We have not found a
collision at any $t\in[0,1/2]$ for the genuine Jacobi
theta transform, and the preceding theorem does not
assert its existence.

\begin{proposition}[An integrable spatial envelope suffices]
\label{prop:logspatialsufficient}
Assume that for every fixed integer $n\ge1$ there is a
nonnegative $a_n\in L^1([0,\tfrac12])$ and a measurable
nonnegative $c_n(s,x)$ such that
\begin{equation}\label{eq:logspatialgate}
\begin{gathered}
k_nU_{n+1}(s,x)\le c_n(s,x)U_n(s,x),\\
0\le c_n(s,x)\le
a_n(s)\bigl(1+\log^2(2+|x|)\bigr)
\end{gathered}
\end{equation}
for all real $x$ and almost every $s\in[0,\tfrac12]$.
Suppose these coefficients have the regularity needed
for the standard parabolic comparison principle
(the genuine Fourier coefficients themselves are smooth
and bounded). Then $U_n(s,x)\ge0$ for all $n\ge1$,
$s\in[0,\tfrac12]$ and real $x$. In particular,
\eqref{eq:logspatialgate} for the Riemann theta heat
family would imply the Riemann hypothesis.
\end{proposition}
\begin{proof}
With $D=\partial_s-\tfrac12\partial_x^2$,
the hierarchy gives
$DU_n=-k_nU_{n+1}\ge-c_nU_n$.
At $s=0$, $H_{1/2}$ has only real zeros, so
$U_n(0,x)\ge0$ by the classical Laguerre--P\'olya
criterion. The linear parabolic problem with
potential $c_n\ge0$ has a positivity-preserving
Feynman--Kac propagator. Under the stated envelope
the stochastic exponential is well-defined on every
finite interval; the Brownian supremum has Gaussian
tails, which control the logarithmic-square spatial
growth. Comparing $U_n$ with the nonnegative solution
started from $U_n(0,\cdot)$ gives $U_n\ge0$.
At $s=1/2$ every generalized Laguerre inequality
holds for $H_0$, so the classical real-zero criterion
implies RH.
\end{proof}

This is a conditional sufficiency theorem, not a proof
of \eqref{eq:logspatialgate}. The collision theorem shows
why a coefficient with a nonintegrable time singularity
is insufficient for the preceding comparison argument.
The genuine theta all-order comparison on $|x|\le6$
established earlier is compatible with this result
but does not discharge the missing global-in-$x$,
all-order estimate. The exact algebraic normal form
is reproduced in
\path{XiKernel/HeatCollisionRate.lean} and
\path{certificates/xi_heat_collision_envelope_exact.py};
the full analytic theta/Feynman--Kac implication is
not yet formalized in Lean.
